% ============================================================
%  Розширення секції «Обговорення» — вставити ПІСЛЯ існуючих
%  трьох підсекцій (після \subsection{Зв'язок з онтологічно"=керованими системами})
%  та ПЕРЕД секцією \section{Обмеження}
% ============================================================


\subsection{Пам'ять як субстрат навчання з підкріпленням}
\label{sec:disc-rlhf-substrate}

Та сама інфраструктура, яка дозволяє одному практику виконувати сотні сесій протягом кількох місяців, \emph{генерує} сигнал переваг, придатний для навчання з підкріпленням на основі зворотного зв'язку від людини (RLHF).
Пам'ять і дані переваг є двома проєкціями одного і того ж довгострокового робочого процесу: кожна корекція після витягування (retrieval-correction edit) є одночасно \emph{збоєм підсистеми пам'яті} --- система не витягнула релевантний контекст --- та \emph{сигналом вирівнювання} --- практик скоригував вихід моделі, породивши пару переваг.

Ця двоїстість не є випадковою: вона становить архітектурну тезу.
Підсистема пам'яті не лише \emph{споживає} дані вирівнювання --- вона їх \emph{продукує}.
Чим інтенсивніша робота практика, тим більше корекцій, тим щільніший сигнал.
Високий операційний темп, що зазвичай ускладнює збір навчальних даних, у цій архітектурі навпаки прискорює їх генерацію.

У бойовому контексті аналогією є зв'язок між інтенсивністю бойових дій та обсягом розвідувальної інформації, що генерується автоматично: кожна корекція оператора під час управління ударним БПЛА --- зміна маршруту, перепризначення цілі, відміна пуску --- є одночасно записом пам'яті місії та зразком для навчання системи на майбутніх аналогічних сценаріях.
Як зазначено у~\citet{ovcharov2026edittrace}, цей двоїстий сигнал може бути використаний для формування датасетів переваг без додаткового навантаження на оператора.


\subsection{Пам'ять як поверхня масштабованого нагляду}
\label{sec:disc-oversight-surface}

Підсистема пам'яті виконує подвійну функцію: обслуговує агента під час виконання місії та генерує сигнал нагляду під час узгодження (reconciliation).
Із зростанням автономності агента --- від коротких інтерактивних сесій до багатогодинних автономних виконань і багатоденних автономних проєктів --- відношення спостережуваних результатів до внутрішніх рішень агента падає.
Нагляд на рівні результатів стає вузькосмуговим каналом, через який повинен пройти весь наглядовий сигнал.

Підсистема пам'яті змінює цю динаміку.
Кожне витягування є \emph{спостережуваною подією} з відомим контекстним вікном, відомим результатом витягування та подальшим сліду редагування.
Узгодження між витягнутим контекстом та поведінкою редагування генерує наглядовий сигнал з гранулярністю окремих редагувань, а не окремих сесій.
Цей сигнал не є безкоштовним --- узгодження вимагає підтвердження практиком кандидатних пар <<витягування--корекція>> --- але маргінальна вартість одного сигналу для практика обмежена часом підтвердження позначеної пари, а не часом оцінювання всієї сесії.

Щільність сигналу на одиницю активності є вищою, ніж при RLHF на рівні результатів: кожна сесія генерує не один бінарний сигнал <<успіх/провал>>, а множину пар <<витягнуто~$X$, скориговано~$Y$>>, де кожна пара є окремим навчальним зразком.

\textbf{Військова аналогія.}
Це аналогічно автоматизованому розбору місії (mission debrief), що генерує дієву розвідувальну інформацію, а не лише журнал місії.
Традиційний розбір --- оператор після місії переказує ключові рішення --- масштабується лінійно з кількістю операторів і місій.
Автоматизований розбір на основі підсистеми пам'яті масштабується з кількістю витягувань, яка зростає пропорційно складності місії.
Чим складніша місія, тим більше сигналу генерується --- саме тоді, коли він найбільш потрібний.


\subsection{Поведінка пам'яті при ескалації подій}
\label{sec:disc-escalation}

Попередні підсекції розглядають підсистему пам'яті в режимі стаціонарного навантаження.
Окремого аналізу потребує поведінка при \emph{ескалації} --- каскадному зростанні кількості одночасних подій, коли інформаційне навантаження на оператора та агента зростає експоненціально.

\textbf{Каскадні сценарії.}
В бойовому контексті ескалація виникає при одночасному обстрілі кількох секторів, радіоелектронному придушенні каналів зв'язку та втраті контакту з частиною рою БПЛА.
У програмному контексті --- при каскадних виробничих інцидентах у кількох сервісах одночасно, коли ротація операторів відбувається посеред багатофронтової аварії.
В обох випадках кількість точок прийняття рішень зростає швидше, ніж пропускна здатність підсистеми пам'яті.

\textbf{Деградований режим.}
Коли частота подій перевищує пропускну здатність обробки, підсистема пам'яті має переключитися з семантичного витягування на \emph{тріаж на основі критичності}.
Семантична подібність перестає бути первинним критерієм ранжування: замість <<що найбільш семантично релевантне>> система відповідає на запитання <<що є найбільш критичним для виживання місії>>.
Формально, функція ранжування змінюється з
\[
  \mathrm{score}(e) = \mathrm{sim}(q, e)
  \quad\longrightarrow\quad
  \mathrm{score}(e) = \alpha \cdot \mathrm{criticality}(e) + (1-\alpha) \cdot \mathrm{sim}(q, e),
\]
де $\alpha \to 1$ при зростанні частоти подій.
Параметр $\alpha$ може визначатися адаптивно --- на основі поточного рівня загрози або частоти вхідних подій за ковзним вікном.

\textbf{Push"=режим під ескалацією.}
В умовах обмеженого токенного бюджету push"=режим стикається з проблемою пріоритизації: які контексти оновлювати першими?
При стаціонарному навантаженні пріоритет визначається часом від останнього оновлення та семантичною дистанцією від поточного стану.
При ескалації пріоритет визначається чергою загроз: контексти, пов'язані з найвищим рівнем загрози, оновлюються першими незалежно від семантичної дистанції.

\textbf{Ротація під час кризи.}
Найгірший сценарій для підсистеми пам'яті --- ротація оператора посеред ескалації.
Дайджест для нового оператора повинен передати не лише статичний стан (<<поточна ситуація>>), а й \emph{динаміку}: що ескалює, що каскадує, які зв'язки між подіями виявлені, які рішення прийняті попереднім оператором і які з них потребують негайного продовження.
Це вимагає розширення формату дайджесту від плоского переліку фактів до темпоральної структури з причинно"=наслідковими зв'язками.

Повна реалізація деградованого режиму виходить за межі поточної роботи (див.~обмеження~(iii) у секції~\ref{sec:limitations}), проте архітектурні точки розширення --- функція ранжування з параметром критичності та черга пріоритетів у push"=режимі --- є частиною проєкту і готові до реалізації.


\subsection{Узагальнення на команди}
\label{sec:disc-teams}

Архітектура у представленому вигляді обслуговує одного практика.
Багатооператорне розгортання породжує три виклики, що не розглядаються у поточній роботі.

\emph{По"=перше}, реєстр принципів (principle ledger) стає багатозаписувачем: кілька операторів можуть одночасно вносити суперечливі принципи або різні формулювання одного й того ж рішення.
Необхідний механізм вирішення конфліктів --- від простого <<останній запис перемагає>> до структурованого голосування або ієрархічного підтвердження старшим оператором.

\emph{По"=друге}, рівень оператора потребує \emph{індивідуальних представлень}: різні оператори мають різні патерни корекцій, різні зони уваги та різний рівень досвіду.
Спільне представлення оператора для всієї команди знищує персоналізацію, що є однією з основних переваг трирівневої декомпозиції.

\emph{По"=третє}, записи пам'яті можуть містити контекст, чутливий до рівня доступу: в бойовому контексті --- рівень секретності, у програмному --- розмежування доступу за проєктами або командами.
Підсистема пам'яті повинна підтримувати контроль доступу на рівні окремих записів.

Трирівнева декомпозиція спроєктована для підтримки цих розширень: предметний рівень є за визначенням спільним, робочий рівень підтримує обмежені представлення через фільтрацію метаданих, а рівень оператора є за природою індивідуальним.
Перехід від одного оператора до команди не потребує перепроєктування архітектури --- лише додавання шарів контролю доступу та узгодження конфліктів поверх існуючих рівнів.


\subsection{Чим це не є}
\label{sec:disc-not}

Архітектура підсистеми пам'яті робочого процесу \emph{не} є заміною моделей з довгим контекстним вікном: вона є витягувальним субстратом, що зменшує обсяг контексту, який повинна обробити модель з довгим вікном.
Моделі на мільйон токенів~\citep{reid2024gemini} дозволяють вмістити більше інформації, але проблема <<загублення в середині>>~\citep{liu2023lost} та деградація на задачах, що вимагають точного витягування з глибоких контекстів~\citep{hsieh2024ruler}, залишаються актуальними.
Підсистема пам'яті не конкурує з довгим контекстом --- вона структурує інформацію так, щоб модель працювала у зоні ефективної уваги.

Це \emph{не} є граф знань: архітектура не містить формальної онтології у класичному розумінні, не використовує RDF"=трійки і не надає SPARQL"=інтерфейс.
Зв'язок з онтологічно"=керованими системами~\citep{palagin2006architecture}, описаний у попередній підсекції, є концептуальною аналогією, а не технічною реалізацією --- рівні пам'яті \emph{зіставляються} з онтологічними концепціями, але не \emph{є} онтологією.

Це \emph{не} є корпоративна система управління знаннями (Enterprise KMS): відсутній документний життєвий цикл, відсутні процеси затвердження, відсутні метадані відповідності (compliance metadata).
Підсистема не призначена для роботи зі структурованими документами, що проходять через формалізовані процеси рецензування.

Це витягувальний субстрат, оптимізований для одного операційного режиму --- довгострокової агентної композиції малою кількістю практиків --- і його проєктні рішення відображають цю специфічність.

\textbf{Військова аналогія.}
Це \emph{не} є повноцінна система C4ISR (командування, управління, зв'язок, комп'ютери, розвідка, спостереження, рекогносцировка).
Це підсистема пам'яті, що інтегрується в існуючу інфраструктуру управління.
Як автомат Калашникова не є армією, а є її складовою частиною, підсистема пам'яті не претендує на замiну системи управління бойовими діями --- вона є компонентом, що забезпечує безперервність інформаційного контексту для тих підсистем, які приймають рішення.
